\documentclass[tikz,border=8pt]{standalone}
\usepackage{ctex}
\usepackage{amsmath}
\usetikzlibrary{calc, arrows.meta, shapes.geometric, positioning}

\begin{document}
\begin{tikzpicture}[
  scale=1.0,
  box/.style={draw, rounded corners=4pt, fill=#1!10, minimum width=3.6cm, minimum height=0.85cm,
              align=center, font=\small},
  arrow/.style={-{Stealth[length=6pt]}, thick, gray},
]

% 标题
\node[font=\large\bfseries] (title) at (4, 6.5) {动点轨迹方程——5 大方法对比};

% 中心问题框
\node[draw, fill=orange!20, rounded corners=6pt, font=\normalsize, minimum width=4cm, minimum height=1cm,
      align=center] (center) at (4, 5.2) {求动点 $P(x,y)$ 的\\轨迹方程};

% 五个方法框（扇形排布）
\node[box=blue] (m1) at (0.5, 3.5) {① 直接法\\列条件方程\\消元化简};
\node[box=green] (m2) at (2.8, 2.2) {② 定义法\\识别椭圆/双曲线/\\抛物线定义};
\node[box=purple] (m3) at (5.2, 2.2) {③ 几何法\\用几何性质\\（中垂线/圆心等）};
\node[box=red] (m4) at (7.3, 3.5) {④ 参数法\\设 $t$ 表达 $x,y$\\消参得方程};
\node[box=cyan!70!black] (m5) at (4, 0.7) {⑤ 相关点法\\$P$ 依赖 $M$，代入\\$M$ 的曲线方程};

% 箭头
\draw[arrow] (center) -- (m1);
\draw[arrow] (center) -- (m2);
\draw[arrow] (center) -- (m3);
\draw[arrow] (center) -- (m4);
\draw[arrow] (center) -- (m5);

% 各方法的适用提示
\node[font=\footnotesize, align=left, blue!70] at (0.5, 2.55)
  {\textit{适用：条件直接转化}};
\node[font=\footnotesize, align=center, green!50!black] at (2.8, 1.3)
  {\textit{适用：焦距和/差已知}};
\node[font=\footnotesize, align=center, purple!70] at (5.2, 1.3)
  {\textit{适用：等距/等长条件}};
\node[font=\footnotesize, align=center, red!70] at (7.3, 2.55)
  {\textit{适用：三角/圆参数化}};
\node[font=\footnotesize, align=center, cyan!60!black] at (4, -0.05)
  {\textit{适用：$P$ 与另一动点 $M$ 相关}};

% 分隔线
\draw[gray!40, very thin] (0, -0.4) -- (8.2, -0.4);
\node[font=\footnotesize, gray] at (4, -0.65) {选方法的关键：读懂题目给的是哪类条件，直接对号入座};

\end{tikzpicture}
\end{document}
